\documentclass[aspectratio=169,11pt]{beamer}
\usepackage[T1]{fontenc}
\usepackage[utf8]{inputenc}
\usepackage[brazil]{babel}
\usepackage{lmodern}
\usepackage{amsmath,amssymb,booktabs,array}
\usepackage{tikz}
\usetikzlibrary{automata,positioning,arrows.meta}
\usepackage{listings}
\definecolor{azul}{RGB}{20,49,78}
\definecolor{verde}{RGB}{0,121,111}
\definecolor{cinza}{RGB}{244,247,250}
\lstset{language=Python,basicstyle=\ttfamily\scriptsize,
 keywordstyle=\color{azul}\bfseries,commentstyle=\color{verde},
 stringstyle=\color{verde},breaklines=true,showstringspaces=false,
 frame=single,rulecolor=\color{azul!20},backgroundcolor=\color{cinza},
 columns=fullflexible,keepspaces=true,tabsize=4}
\newcommand{\q}[1]{q_{\mathrm{#1}}}
\newcommand{\MiniL}{\textsc{MiniL}}
\newcommand{\grupo}{Gabriel Holanda, Gabriel Torres, Victor Milito,\\Paulo Júnior, Vinicius Kaua e Pablo Café}
\tikzset{every state/.style={draw=azul,thick,minimum size=8mm,font=\small},
 every edge/.style={draw=azul,>=Stealth},every edge quotes/.style={font=\scriptsize}}

\usetheme{default}
\useinnertheme{rectangles}
\setbeamercolor{frametitle}{bg=azul,fg=white}
\setbeamercolor{block title}{bg=azul,fg=white}
\setbeamercolor{block body}{bg=cinza,fg=black}
\setbeamercolor{structure}{fg=azul}
\setbeamercolor{palette primary}{bg=azul,fg=white}
\setbeamercolor{palette secondary}{bg=verde,fg=white}
\setbeamercolor{palette tertiary}{bg=azul,fg=white}
\setbeamertemplate{navigation symbols}{}
\setbeamertemplate{footline}[frame number]
\title[MiniL]{O Compilador Mínimo}
\subtitle{Modelando um reconhecedor léxico para a MiniL}
\author[Equipe]{\small\grupo}
\institute{Linguagens Formais e Autômatos Finitos}
\date{Revisão 1.1.0 --- 13/11/2026}
\begin{document}
\begin{frame}\titlepage\end{frame}

\begin{frame}{1. O que vamos reconhecer?}
\begin{block}{Entrada}
\texttt{let total = 12.50;}
\end{block}
\begin{center}
\begin{tabular}{ccccc}
\toprule \texttt{let}&\texttt{total}&\texttt{=}&\texttt{12.50}&\texttt{;}\\
KW\_LET & ID & OP & FLOAT & SEP\\\bottomrule
\end{tabular}
\end{center}
\begin{itemize}
\item Lexema: texto encontrado. Token: categoria e lexema, com posição.
\item O scanner é a etapa léxica de um compilador.
\item Não verificamos sintaxe, tipos ou execução do programa.
\end{itemize}
\end{frame}

\begin{frame}{2. Um contrato pequeno e explícito}
\begin{columns}[T]
\column{.53\textwidth}
\begin{block}{Reconhecemos}
\begin{itemize}
\item Reservadas: \texttt{let, if, else, print}.
\item ID: letra ASCII ou \texttt{\_}; depois, também dígitos.
\item Inteiros e decimais com ponto.
\item Operadores: \texttt{+ - * / =}.
\item Separadores: \texttt{( ) ;}.
\end{itemize}
\end{block}
\column{.43\textwidth}
\begin{exampleblock}{Decisões importantes}
\texttt{12.5}: FLOAT.\\[2mm]
\texttt{12.}: erro no ponto.\\[2mm]
\texttt{-3}: OP e INT.\\[2mm]
\texttt{ifx}: um ID.\\[2mm]
Espaços são descartados.
\end{exampleblock}
\end{columns}
\end{frame}

\begin{frame}{3. Expressões regulares: a forma dos lexemas}
\begin{align*}
L&=[\texttt{A-Z a-z \_}],\qquad D=[\texttt{0-9}]\\
\mathrm{ID}&=L(L\mid D)^*\\
\mathrm{INT}&=D^+\\
\mathrm{FLOAT}&=D^+\,\texttt{.}\,D^+
\end{align*}
\begin{itemize}
\item OP e SEP: um caractere das respectivas classes $O$ e $S$.
\item WS: $W^+$, com espaço, tabulação, retorno de carro e quebra de linha.
\item Reservadas: consulta à tabela após consumir o ID completo.
\end{itemize}
\end{frame}

\begin{frame}{4. A definição formal}
\[
M=(Q,\Sigma,\delta,\q{0},F)
\]
\begin{align*}
Q&=\{\q{0},\q{I},\q{N},\q{Dot},\q{F},\q{Op},\q{Sep},\q{W},\q{Err}\},\\
F&=\{\q{I},\q{N},\q{F},\q{Op},\q{Sep},\q{W}\},\\
\delta&:Q\times\Sigma\to Q.
\end{align*}
\begin{itemize}
\item $\Sigma$: conjunto finito de valores de caracteres; classes
$L,D,\{.\},O,S,W,X$, sendo $X$ o restante.
\item $\q{0}$: inicial. $\q{Err}$: erro absorvente.
\item Uma palavra $w$ é aceita quando $\delta^*(\q{0},w)\in F$.
\end{itemize}
\end{frame}

\begin{frame}{5. O AFD original: nove estados}
\begin{center}\resizebox{!}{.64\textheight}{\begin{tikzpicture}[->,auto,node distance=17mm and 20mm,scale=.92,transform shape]
\node[state,initial] (s) {$\q{0}$};
\node[state,accepting,above right=of s] (i) {$\q{I}$};
\node[state,accepting,right=of s] (n) {$\q{N}$};
\node[state,right=of n] (d) {$\q{Dot}$};
\node[state,accepting,right=of d] (f) {$\q{F}$};
\node[state,accepting,below right=of s] (o) {$\q{Op}$};
\node[state,accepting,right=of o] (p) {$\q{Sep}$};
\node[state,accepting,below=of s] (w) {$\q{W}$};
\node[state,right=of p] (e) {$\q{Err}$};
\path (s) edge node {$L$} (i)
 (i) edge[loop above] node {$L,D$} (i)
 (s) edge node {$D$} (n)
 (n) edge[loop above] node {$D$} (n)
 (n) edge node {$.$} (d)
 (d) edge node {$D$} (f)
 (f) edge[loop above] node {$D$} (f)
 (s) edge node[left] {$O$} (o)
 (s) edge[bend left=12] node[above,pos=.75] {$S$} (p)
 (s) edge node[left] {$W$} (w)
 (w) edge[loop left] node {$W$} (w)
 (e) edge[loop right] node {$\Sigma$} (e);
\end{tikzpicture}
}\end{center}
\small Arestas omitidas vão para $\q{Err}$. Círculos duplos são finais.
O laço em $\q{Err}$ abrange $\Sigma$. A tabela completa está no relatório.
\end{frame}

\begin{frame}{6. Por que o ponto precisa de um estado?}
\begin{block}{Rastreamento de \texttt{12.50}}
\[
\q{0}\xrightarrow{1}\q{N}\xrightarrow{2}\q{N}
\xrightarrow{.}\q{Dot}\xrightarrow{5}\q{F}\xrightarrow{0}\q{F}
\]
\end{block}
\begin{itemize}
\item $\q{N}$ já aceita inteiro.
\item $\q{Dot}$ ainda não aceita: falta um dígito após o ponto.
\item $\q{F}$ aceita decimal e repete dígitos.
\end{itemize}
\begin{alertblock}{Caso negativo}
\texttt{12.} termina em $\q{Dot}\notin F$: não é um lexema válido.
\end{alertblock}
\end{frame}

\begin{frame}{7. Minimização: primeiras separações}
\small
Todos os estados são alcançáveis: $\varepsilon$, \texttt{a}, \texttt{0},
\texttt{0.}, \texttt{0.0}, \texttt{+}, \texttt{;}, espaço e \texttt{@}.
\begin{align*}
\mathcal{P}_0=\{&\{\q{I},\q{N},\q{F},\q{Op},\q{Sep},\q{W}\},
\{\q{0},\q{Dot},\q{Err}\}\}.\\[2mm]
\mathcal{P}_1=\{&\{\q{I}\},\{\q{N},\q{F}\},\{\q{Op},\q{Sep}\},
\{\q{W}\},\{\q{0}\},\{\q{Dot}\},\{\q{Err}\}\}.
\end{align*}
\begin{block}{Regra de refinamento de Moore}
Dentro de cada bloco, comparamos os blocos de destino para cada classe.
Destinos em blocos diferentes obrigam a separação.
\end{block}
$\q{N}$ e $\q{F}$ ainda ficam juntos na primeira rodada.
\end{frame}

\begin{frame}{8. Minimização: estabilidade e resultado}
\[
\delta(\q{N},\texttt{.})=\q{Dot},\qquad
\delta(\q{F},\texttt{.})=\q{Err}.
\]
Agora esses destinos pertencem a blocos diferentes. Portanto:
\begin{align*}
\mathcal{P}_2=\{&\{\q{0}\},\{\q{I}\},\{\q{N}\},\{\q{Dot}\},\\
&\{\q{F}\},\{\q{Op},\q{Sep}\},\{\q{W}\},\{\q{Err}\}\}.
\end{align*}
\begin{block}{Resultado}
$\mathcal{P}_3=\mathcal{P}_2$: oito estados. Só OP e SEP se fundem.
Todos os blocos são alcançáveis e distinguíveis por sufixos.
\end{block}
Exemplo: o sufixo \texttt{.0} distingue $\q{N}$ de $\q{F}$.
\end{frame}

\begin{frame}{9. O reconhecedor mínimo e as categorias}
\begin{columns}
\column{.66\textwidth}
\resizebox{.97\linewidth}{!}{\begin{tikzpicture}[->,auto,node distance=17mm and 20mm,scale=.92,transform shape]
\node[state,initial] (s) {$[\q{0}]$};
\node[state,accepting,above right=of s] (i) {$[\q{I}]$};
\node[state,accepting,right=of s] (n) {$[\q{N}]$};
\node[state,right=of n] (d) {$[\q{Dot}]$};
\node[state,accepting,right=of d] (f) {$[\q{F}]$};
\node[state,accepting,below right=of s] (t) {$\q{T}$};
\node[state,accepting,below=of s] (w) {$[\q{W}]$};
\node[state,right=of t] (e) {$[\q{Err}]$};
\path (s) edge node {$L$} (i)
 (i) edge[loop above] node {$L,D$} (i)
 (s) edge node {$D$} (n)
 (n) edge[loop above] node {$D$} (n)
 (n) edge node {$.$} (d)
 (d) edge node {$D$} (f)
 (f) edge[loop above] node {$D$} (f)
 (s) edge node[right] {$O,S$} (t)
 (s) edge node[left] {$W$} (w)
 (w) edge[loop left] node {$W$} (w)
 (e) edge[loop right] node {$\Sigma$} (e);
\end{tikzpicture}
}
\column{.32\textwidth}
\small
\begin{block}{$\q{T}=\{\q{Op},\q{Sep}\}$}
Os dois estados aceitam e depois sempre vão para erro.
\end{block}
\textbf{Aceitar não é classificar.}\\[2mm]
Após a fusão, OP/SEP é recuperado pelo lexema.\\[2mm]
O scanner fornecido usa o AFD original; o mínimo é calculado e verificado.
\end{columns}
\vspace{1mm}\scriptsize Transições omitidas levam ao estado de erro.
\end{frame}

\begin{frame}{10. Como separar um arquivo em tokens?}
\begin{enumerate}
\item Reiniciar em $\q{0}$ na posição atual.
\item Avançar e guardar a última posição em estado final.
\item Ao parar, retornar à última aceitação e classificar o lexema.
\item Descartar espaços; emitir os demais tokens e repetir.
\end{enumerate}
\begin{exampleblock}{Maior prefixo aceito}
\texttt{ifx} vira um ID inteiro.\\
\texttt{12.} reconhece \texttt{12}, retorna ao ponto e informa erro.\\
\texttt{a+b} vira três tokens, embora não seja um único lexema.
\end{exampleblock}
Neste AFD, o retrocesso é limitado; o scanner custa $O(n)$.
\end{frame}

\begin{frame}[fragile]{11. Implementação: a última aceitação}
\begin{lstlisting}
q, j, last = 'q0', i, None
while j < len(source):
    q = DELTA[q][char_class(source[j])]
    if q == DEAD:
        break
    j += 1
    if q in FINAL:
        last = (j, q)
if last is None:
    raise LexicalError(...)
end, accepted = last
lexeme = source[i:end]
\end{lstlisting}
\small Trecho resumido de \texttt{tokenize()}; a exceção completa está no arquivo.
\texttt{DELTA} é a própria tabela formal. Nenhuma regex é executada pelo scanner.
\end{frame}

\begin{frame}[fragile]{12. Demonstração ao vivo e evidências}
\begin{lstlisting}[language={}]
python codigo/scanner.py exemplos/exemplo.min --minimizar
python -m unittest discover -s testes -v
python codigo/scanner.py exemplos/invalido.min
\end{lstlisting}
\begin{itemize}
\item Observar tokens, categorias e linha:coluna.
\item Conferir nove estados no original e oito no quociente.
\item 18 testes: fronteiras, erros, posições e propriedades formais.
\item Equivalência: explorar o produto alcançável dos dois AFDs.
\end{itemize}
\begin{alertblock}{Erro esperado na demonstração negativa}
\texttt{let total = 12.;}\quad acusa o ponto em \texttt{1:15}.
\end{alertblock}
\end{frame}

\begin{frame}{13. O que demonstramos}
\begin{itemize}
\item Expressões regulares definem os lexemas da MiniL.
\item Um AFD completo reconhece exatamente essas formas.
\item Moore reduz nove estados a oito, com prova de minimalidade.
\item O scanner aplica maior prefixo e preserva a posição dos tokens.
\item Versão 1.1.0: código tipado, JSON e 18 testes; revisão de 13/11/2026.
\end{itemize}
\begin{block}{Limite do projeto}
Reconhecimento léxico é uma etapa. Verificar estrutura e significado exigiria
análises sintática e semântica adicionais.
\end{block}
\end{frame}

\begin{frame}{Referências}
\small
SIPSER, Michael. \emph{Introduction to the Theory of Computation}.
3. ed. Cengage Learning, 2013.\\[5mm]
AHO, Alfred V.; LAM, Monica S.; SETHI, Ravi; ULLMAN, Jeffrey D.
\emph{Compilers: Principles, Techniques, and Tools}.
2. ed. Addison-Wesley, 2006.\\[3mm]
Python Software Foundation. \emph{Python 3.14 Documentation}.\\
\url{https://docs.python.org/3.14/}. Acesso em: 13/11/2026.\\[4mm]
\textbf{Materiais da entrega:} relatório article, apresentação beamer,
relatório para o cliente, diagramas TikZ, scanner Python e testes.
\end{frame}
\end{document}
